\chapter{Capstone project}
\label{ch:capstone}

\begin{quote}
\emph{``The best way to learn data preprocessing is to preprocess data.''}
\end{quote}

% ============================================================
\section{Capstone overview}

In this final chapter, you will apply everything from Chapters~1--9 to a complete, end-to-end preprocessing project. You will choose one of five project options, each using a real, messy dataset that requires substantial cleaning and transformation before analysis.

\begin{datatip}
Each project is designed to take approximately 3 hours. You are expected to deliver: (1)~a Jupyter notebook with documented code, (2)~a clean, analysis-ready dataset saved as Parquet, and (3)~a brief written report (1--2 pages) summarizing your preprocessing decisions and their justification.
\end{datatip}

% ============================================================
\section{Evaluation criteria}

Your project will be evaluated on five criteria:

\begin{longtable}{p{4cm}p{3cm}p{6cm}}
\toprule
\textbf{Criterion} & \textbf{Weight} & \textbf{Description} \\
\midrule
Data quality report & 15\% & Thorough initial assessment of data quality \\
Missing data strategy & 20\% & Justified choice of imputation method(s) \\
Outlier handling & 15\% & Detection, investigation, and treatment \\
Feature engineering & 25\% & Creativity and domain relevance of new features \\
Pipeline and reproducibility & 25\% & Complete, leakage-free, saved pipeline \\
\bottomrule
\end{longtable}

% ============================================================
\section{Project 1: Melbourne Housing price prediction}

\textbf{Dataset:} Melbourne Housing Snapshot (Kaggle)\\
\textbf{Goal:} Prepare the dataset for a house price regression model.

\subsection{Requirements}

\begin{minted}{python}
import pandas as pd

# Load the raw dataset
df = pd.read_csv("melb_data.csv")
print(f"Shape: {df.shape}")
print(f"Missing values:\n{df.isnull().sum()}")
\end{minted}

Tasks:
\begin{enumerate}
  \item Produce a comprehensive data quality report (Chapter~1).
  \item Handle missing values in \texttt{BuildingArea}, \texttt{YearBuilt}, \texttt{Car}, \texttt{CouncilArea} using at least two different imputation strategies. Justify your choices (Chapter~3).
  \item Detect and treat outliers in \texttt{Price}, \texttt{Landsize}, and \texttt{BuildingArea} (Chapter~4).
  \item Encode categorical variables (\texttt{Type}, \texttt{Method}, \texttt{Regionname}, \texttt{Suburb}). For high-cardinality \texttt{Suburb}, use target encoding or frequency encoding (Chapter~5).
  \item Engineer at least 5 new features: price per room, building age, distance categories, suburb-level averages, seasonal sale indicator (Chapter~6).
  \item Build a complete \texttt{Pipeline} with \texttt{ColumnTransformer}. Save with \texttt{joblib} (Chapter~9).
  \item Train a baseline model and report $R^2$ on a held-out test set.
\end{enumerate}

% ============================================================
\section{Project 2: Titanic survival classification}

\textbf{Dataset:} Titanic (Kaggle)\\
\textbf{Goal:} Prepare the dataset for a survival classification model.

\subsection{Requirements}

\begin{minted}{python}
url = ("https://raw.githubusercontent.com/datasciencedojo/"
       "datasets/master/titanic.csv")
df = pd.read_csv(url)
print(f"Shape: {df.shape}")
print(f"Survival rate: {df['Survived'].mean():.2%}")
\end{minted}

Tasks:
\begin{enumerate}
  \item Analyze missing data patterns. Investigate whether \texttt{Age} missingness is MCAR or MAR by comparing it against \texttt{Pclass} and \texttt{Sex} (Chapter~3).
  \item Impute \texttt{Age} using KNN imputation. Compare with median imputation by evaluating downstream model accuracy.
  \item Extract title from \texttt{Name} (Mr, Mrs, Miss, Master, etc.) and encode it (Chapters~5 and 7).
  \item Create features: \texttt{FamilySize} = \texttt{SibSp} + \texttt{Parch} + 1, \texttt{IsAlone}, \texttt{FarePerPerson}, cabin deck letter (Chapter~6).
  \item Build a complete pipeline with cross-validated evaluation.
\end{enumerate}

% ============================================================
\section{Project 3: FIFA 21 player value estimation}

\textbf{Dataset:} FIFA 21 Raw Data (Kaggle)\\
\textbf{Goal:} Clean the notoriously messy FIFA dataset and prepare it for player value prediction.

\subsection{Requirements}

\begin{minted}{python}
df = pd.read_csv("fifa21_raw.csv")
print(f"Shape: {df.shape}")
# Note: many columns have string-encoded numbers (e.g., "€110M", "190lbs")
print(df[["Value", "Wage", "Weight", "Height"]].head())
\end{minted}

Tasks:
\begin{enumerate}
  \item Parse string-encoded currency values (\texttt{Value}, \texttt{Wage}): convert ``\texteuro110M'' to 110000000, ``\texteuro220K'' to 220000.
  \item Parse height (``5'11'') and weight (``190lbs'') to metric units.
  \item Handle the \texttt{Hits} and \texttt{W/F}, \texttt{SM}, \texttt{IR} columns (star ratings stored as strings).
  \item Detect and treat outliers in player values (many free agents have value 0).
  \item Engineer position-based features, age-based potential features, and physical attribute ratios.
  \item Build a clean pipeline from raw CSV to model-ready features.
\end{enumerate}

% ============================================================
\section{Project 4: Jena Climate time series forecasting}

\textbf{Dataset:} Jena Climate 2009--2016 (Kaggle)\\
\textbf{Goal:} Prepare the dataset for temperature forecasting.

\subsection{Requirements}

\begin{minted}{python}
df = pd.read_csv("jena_climate_2009_2016.csv")
df["datetime"] = pd.to_datetime(df["Date Time"],
                                 format="%d.%m.%Y %H:%M:%S")
df = df.set_index("datetime")
print(f"Shape: {df.shape}")
print(f"Date range: {df.index.min()} to {df.index.max()}")
\end{minted}

Tasks:
\begin{enumerate}
  \item Resample from 10-minute to hourly and daily frequencies (Chapter~8).
  \item Detect and handle sensor anomalies (sudden spikes, constant readings).
  \item Create lag features: 1h, 6h, 12h, 24h, 7d for temperature.
  \item Create rolling statistics: 24h mean, 7d mean, 24h std.
  \item Decompose into trend, seasonality, and residual. Use each as a feature.
  \item Add cyclical encoding for hour of day, day of week, and month.
  \item Build a pipeline and train a baseline model for 24-hour-ahead temperature prediction.
\end{enumerate}

% ============================================================
\section{Project 5: 20 Newsgroups text classification}

\textbf{Dataset:} 20 Newsgroups (scikit-learn)\\
\textbf{Goal:} Preprocess text data for topic classification.

\subsection{Requirements}

\begin{minted}{python}
from sklearn.datasets import fetch_20newsgroups

# Use a subset of 5 categories
categories = ["sci.space", "rec.sport.baseball", "comp.graphics",
              "talk.politics.mideast", "soc.religion.christian"]
newsgroups = fetch_20newsgroups(subset="all", categories=categories,
                                remove=("headers", "footers", "quotes"))
print(f"Documents: {len(newsgroups.data)}")
print(f"Categories: {newsgroups.target_names}")
\end{minted}

Tasks:
\begin{enumerate}
  \item Build a text cleaning pipeline: lowercase, remove emails/URLs/numbers, remove special characters (Chapter~7).
  \item Tokenize, remove stopwords, and lemmatize.
  \item Create a TF-IDF matrix with appropriate parameters (\texttt{min\_df}, \texttt{max\_df}, \texttt{ngram\_range}).
  \item Add document-level features: document length, average word length, vocabulary richness (unique words / total words).
  \item Build a pipeline combining text preprocessing and a classifier.
  \item Evaluate with cross-validation and a confusion matrix.
\end{enumerate}

% ============================================================
\section{Deliverables checklist}

\begin{minted}{python}
# Template for the final deliverable check
deliverables = {
    "data_quality_report": False,    # Section with quality metrics
    "missing_data_handled": False,    # Justified imputation
    "outliers_handled": False,        # Detection + treatment
    "features_engineered": False,     # >= 5 new features
    "pipeline_built": False,          # sklearn Pipeline
    "pipeline_saved": False,          # joblib file
    "baseline_model": False,          # Train + evaluate
    "report_written": False,          # 1-2 page summary
}

# Check off each item as you complete it
for item, done in deliverables.items():
    status = "DONE" if done else "TODO"
    print(f"  [{status}] {item}")
\end{minted}

% ============================================================
\section{Exercises}

\begin{exercise}
\begin{enumerate}
  \item Choose one of the five projects above. Complete all required tasks and submit your Jupyter notebook, saved pipeline, and written report.
  \item (Bonus) After completing your main project, apply the same preprocessing approach to a second project. Compare the challenges and discuss how preprocessing strategies differ across domains.
  \item (Bonus) Deploy your saved pipeline as a simple prediction function: write a script that loads the pipeline from \texttt{joblib} and accepts new data from the command line or a CSV file.
\end{enumerate}
\end{exercise}

% ============================================================
\section{Chapter summary}

\begin{itemize}
  \item A capstone project integrates all preprocessing skills: loading, cleaning, imputation, outlier handling, encoding, scaling, feature engineering, and pipeline automation.
  \item Real datasets require domain-specific decisions at every step --- there is no single ``correct'' preprocessing recipe.
  \item Reproducibility (via pipelines and saved artifacts) is as important as accuracy.
  \item Documentation of preprocessing choices and their justification is a professional requirement.
  \item The five project options cover tabular regression, tabular classification, messy data cleaning, time series, and text --- the main data preprocessing domains.
\end{itemize}
